\section{Return measurement}

\begin{qfdefinition}[Return measures: holding-period, price and total returns; time-weighted and money-weighted returns; cumulative, arithmetic-mean, geometric-mean and annualised returns; log returns; gross and net, real and leveraged returns]\label{def:return-measures}
\qffield{Statement} GIPS 2020 defines a \emph{time-weighted return} (TWR) as ``a method of calculating period-by-period returns that reflects the change in value and negates the effects of external cash flows'', and a \emph{money-weighted return} (MWR) as ``the return for a period that reflects the change in value and the timing and size of external cash flows'' \citep[Glossary, pp.~83 and~78]{wide-gips-2020}. An \emph{external cash flow} is capital entering or leaving a portfolio; dividend and interest income are not external cash flows (p.~75). For portfolios in composites, firms must value at least monthly and on the date of all large cash flows, calculate sub-period returns at large cash flows if daily returns are not calculated, with exceptions for private-market investment portfolios, and ``geometrically link periodic and sub-period returns'' (provisions 2.A.23--2.A.24, p.~11).

When calculating money-weighted returns, firms must ``calculate annualized since-inception money-weighted returns'' using daily external cash flows (2.A.29, p.~12), but ``returns for periods of less than one year must not be annualized'' (2.A.12, p.~9); accordingly a composite money-weighted report presents the annualised since-inception return through the most recent annual period end, and, for a track record shorter than a full year, the \emph{non-annualised} since-inception return through the initial annual period end (5.A.1.a--b, p.~33). A \emph{gross-of-fees} return is ``reduced by any transaction costs'', and a \emph{net-of-fees} return is ``the gross-of-fees return reduced by investment management fees'' (Glossary, pp.~76 and~78).

A share's \emph{holding-period return} adds income to the price change: total dollar return is dividend income plus capital gain, and total percent return is dividend yield plus capital-gain yield, each divided by the purchase price; the \emph{price return} is the capital-gain part alone \citep[sec.~15.1]{principles-of-finance-2026}. Indexes are published in both forms: in the FT Wilshire 5000 series ``the reinvestment of all dividend payments is reflected in the total return indexes at the close on the ex-dividend date but is not reflected in the price indexes'' \citep[sec.~1.1, p.~3]{ft-wilshire-5000-2026}. Over several periods the \emph{arithmetic} mean adds the returns and divides by their number; the \emph{geometric} mean multiplies the growth factors $1+r_i$, takes the $n$th root and subtracts one. The geometric average ``tells you what you actually earned per year on average, compounded annually'', the arithmetic ``what you earned in a typical year''; the geometric is smaller unless all the returns are equal \citep[sec.~15.1]{principles-of-finance-2026}.

\qffield{Notation} One period from price $P_0$ to $P_1$ with income $D_1$ received at the end: price return $P_1/P_0-1$, income return $D_1/P_0$, total return $(P_1+D_1)/P_0-1$, their sum. A total-return index multiplies its level each period by $(P_t+D_t)/P_{t-1}$, which is the return of holding the share and reinvesting the dividend at the ex-dividend price. Arithmetic mean $\bar r=\frac1n\sum r_i$; geometric mean $g=\bigl(\prod(1+r_i)\bigr)^{1/n}-1$; $g\le\bar r$, with equality only when all $r_i$ are equal (the inequality of arithmetic and geometric means). For small returns $g\approx\bar r-\sigma^2/2$, with $\sigma^2$ the population variance of the $r_i$; this is an approximation, not an identity.

Interpretation, derived here: with sub-period returns $r_1,\dots,r_n$, each measured between external cash flows, the cumulative TWR is $R=\prod_{i=1}^n(1+r_i)-1$; the geometric mean per sub-period, for equal sub-periods, is $(1+R)^{1/n}-1$; and a return $R$ over $d$ days annualises mathematically to $(1+R)^{365/d}-1$. That formula is arithmetic; whether an annualised figure may be \emph{presented} is a separate rule, and under GIPS it may not be for periods shorter than a year (2.A.12). The MWR is the internal rate of return $x$ solving $\sum_t CF_t(1+x)^{-t}=0$, with contributions negative and withdrawals and the final value positive, from the investor's side. A continuously compounded (log) return is $\ln(1+R)$. A real return is $(1+r_{\mathrm{nom}})/(1+\pi)-1$ for inflation $\pi$.

With own capital $V_0$ and borrowing $V_B$ at rate $r_B$ invested at return $r$ for one period, the return on own capital is
\[ r_L=\frac{r(V_0+V_B)-r_BV_B}{V_0}=r+\frac{V_B}{V_0}(r-r_B). \]

\qffield{Conditions} The TWR removes the effect of the investor's timing of contributions and withdrawals, which suits judging a manager who does not control them; the MWR includes it, and GIPS allows firms to present it only when they control the external cash flows and the composite or pooled fund is closed-end, fixed-life, fixed-commitment or significantly illiquid (1.A.35, pp.~5--6). The IRR need not exist or be unique: the polynomial in $(1+x)^{-1}$ can have several positive roots when cash flows change sign more than once, or none. Flows of $-1$, $+2.5$, $-1.54$ have IRRs of 10 and 40 percent; flows of $-1$, $+1$, $-1$ have none, since their net present value is at most $-0.75$.

Annualising a short-period return assumes it would recur; GIPS forbids annualising returns for periods of less than one year (2.A.12, a calculation-methodology provision that also governs what is presented), and gives no reason; interpretation, that assumption is the obvious one. For periods of a year or more, annualising is a change of units. The leveraged formula assumes one period, borrowing at a fixed rate $r_B$ with no margin calls, fees or forced sales, and $r$ and $r_B$ measured over the same period.

\qffield{Practice} Examples, original. A share bought at 50.00 pays a dividend of 1.20 and is sold at 54.00 a year later: price return 8.0 percent, income return 2.4 percent, total return 10.4 percent. A share at 100 goes ex-dividend with a dividend of 2 at a price of 102 and ends at 105: the price index rises 5.0 percent, the total-return index 7.0588 percent, the return on one share whose dividend buys $2/102$ more shares. Returns of $+20$, $-10$ and $+5$ percent have an arithmetic mean of 5.0 percent and a geometric mean of 4.2808 percent, which compounds to the realised growth factor 1.134; the approximation $\bar r-\sigma^2/2$ gives 4.25 percent. A portfolio of 1{,}000 that earns 8.0 percent after transaction costs (gross of fees) and pays a management fee of 1.0 percent of opening assets at the period end returns $(1{,}080-10)/1{,}000-1=7.0$ percent net of fees.

Example: \$100 grows 10 percent to \$110, \$100 is added, and the \$210 falls 10 percent to \$189. The TWR is $1.10\times0.90-1=-1.0$ percent \emph{cumulative over both periods}, unaffected by the addition, a geometric mean of $-0.50$ percent per sub-period. The MWR solves $-100-100(1+x)^{-1}+189(1+x)^{-2}=0$, giving $x=-3.71$ percent \emph{per sub-period}, or $-7.29$ percent compounded over both. On matching horizons the MWR is lower, $-7.29$ against $-1.0$ percent over the two periods and $-3.71$ against $-0.50$ percent per sub-period, because more money was invested in the losing period. Sub-period returns of $+10$, $-10$ and $+5$ percent compound to 3.95 percent, a geometric mean of 1.30 percent per period.

Four percent over 200 days annualises mathematically to 7.42 percent, a figure GIPS would not allow a firm to present (2.A.12). With 8 percent nominal and 3 percent inflation the real return is 4.85 percent, not 5. With $V_0=100$, $V_B=50$, $r=10$ and $r_B=4$ percent, $r_L=13$ percent. Harris's IRR and holding-period return, and the effect of non-pro-rata flows, are in \hyperref[meth:relative-performance-evaluation]{relative performance evaluation}.

\qffield{Pitfalls} The cumulative TWR and the geometric mean per sub-period are different quantities (Notation). Use the geometric mean to describe realised compound growth and the arithmetic mean as the estimate of a single period's expected return; quoting an arithmetic mean as an annual compound rate overstates growth. A price index omits dividends, so its return understates the holder's return; a total-return index assumes the dividend is reinvested at the ex-dividend date. Sub-periods must be split at external cash flows. A real return adjusts for inflation alone. After-tax returns depend on jurisdiction, investor and realisation timing. The return identities are algebraic derivations.

\qffield{See also} \hyperref[meth:relative-performance-evaluation]{relative performance evaluation}; \hyperref[def:interest-rate-measures]{interest-rate measures} (compounding and annualising rates).

\qffield{Further reading} Composite construction and disclosures under GIPS go beyond these definitions.

\end{qfdefinition}
